\documentclass[../../main.tex]{subfiles}
\begin{document}

{\bf [解答]}

(1) を (2) (3) に代入して
\begin{align}
  & -f''(x) \ge f(x) \label{eq:1} \\
  & f'(0) = 0  \label{eq:2} \\
  & 0 \le x \le a \implies f(x) > 0  \label{eq:3}
\end{align}
である．

\paragraph{(i)}

$0 \le t \le a$ の時\cref{eq:1,eq:3}から
\begin{align}
  -f''(t) > 0 \quad\therefore f''(t) < 0
\end{align}
だから，これを$[0, x]\ (0 < x \le a)$ で積分して
\begin{align}
  & 0 > \int_{0}^{x}f''(t)dt = f'(x) - f'(0) =  f'(x)  \quad (\because \text{\cref{eq:2}}) \\
  \therefore
  & f'(x) < 0  \label{eq:4}
\end{align}
だから $x = a$ として(1)の条件から$f'(a)=-g(a)$だから
\begin{align}
  f'(a) < 0 \iff 0 < g(a)
\end{align}
となり題意は示された．

\paragraph{(ii)}

\begin{figure}[htb]
  \centering
  \begin{tikzpicture}[x=2.2cm, y=1.8cm, >=stealth,
    every node/.style={font=\small}]
    \coordinate (O) at (0, 0);
    \coordinate (A) at (1.0, 0);
    \coordinate (P) at (1.0, 1.5);
    \coordinate (Y0) at (2.0, 0);
    \coordinate (XT) at (2.5, 0);

    % 座標軸
    \draw[->, thick, gray!80!black] (-0.3, 0) -- (3.0, 0) node[right] {$x$};
    \draw[->, thick, gray!80!black] (0, -0.4) -- (0, 2.4) node[above] {$y$};
    \node[below left=2pt] at (O) {$O$};

    % 垂線・補助線
    \draw[dashed, gray!70] (A) -- (P);
    \draw[dashed, gray!70] (0, 1.5) -- (P);
    \node[left=2pt, font=\footnotesize] at (0, 1.5) {$f(a)$};

    % 曲線 y = f(x)
    \draw[line width=1.3pt, blue!85!black] plot[domain=0:2.15, samples=50] (\x, {2.0 - 0.5*\x*\x});
    \node[above=2pt, blue!85!black] at (0.3, 1.95) {$y = f(x)$};

    % 接線
    \draw[line width=1.0pt, teal!80!black] plot[domain=0.4:2.7, samples=10] (\x, {2.5 - \x});
    \node[above right=2pt, teal!85!black, font=\footnotesize] at (1.8, 0.7) {接線};

    % 主要点
    \node[circle, fill=blue!85!black, inner sep=1.4pt] at (P) {};
    \node[above right=2pt] at (P) {$(a, f(a))$};

    \node[circle, fill=blue!85!black, inner sep=1.4pt] at (A) {};
    \node[below=3pt] at (A) {$a$};

    \node[circle, fill=red!85!black, inner sep=1.4pt] at (Y0) {};
    \node[below=3pt, red!85!black] at (Y0) {$y$};

    \node[circle, fill=teal!80!black, inner sep=1.4pt] at (XT) {};
    \node[below=3pt, teal!85!black, font=\footnotesize] at (XT) {$a + \frac{f(a)}{g(a)}$};
  \end{tikzpicture}
  \caption{$x=a$ における接線と零点 $y$ の位置関係．}
  \label{fig:1977_6_tangent}
\end{figure}

題意を背理法で示す．(3)より$f(a)>0$だから，
$a \le y \le a - \dfrac{f(a)}{f'(a)}$ で常に $f(y) > 0$ と仮定する．
\cref{eq:3}から$0 \le y \le a$でも$f(y)>0$だから，
\begin{align}
  0 \le y \le a - \dfrac{f(a)}{f'(a)} \label{eq:5}
\end{align}
で常に$f(y)>0$となる．よって\cref{eq:1}の区間では
\begin{align}
  f''(y) < 0
\end{align}
となり，$f'(y)$は単調減少である．これと\cref{eq:2}より
\begin{align}
  f'(y) \le 0 \label{eq:6}
\end{align}
となり，$f(y)$は単調減少である．

ここで $f(x)$ には平均値の定理が適用できて，
\begin{align}
  f\left(a - \dfrac{f(a)}{f'(a)}\right) - f(a) = -\dfrac{f(a)}{f'(a)} f'(c) \quad
  \left(a < c < a - \dfrac{f(a)}{f'(a)}\right)
\end{align}
なる $c$ が存在する．変形して
\begin{align}
  f\left(a - \dfrac{f(a)}{f'(a)}\right) = \dfrac{f(a)}{f'(a)} \{ f'(a) - f'(c) \} \label{eq:7}
\end{align}
である．右辺について，$\dfrac{f(a)}{f'(a)}<0$と，$f'(x)$が単調減少だから$f'(a)>f'(c)$だから右辺は負である．一方で左辺は仮定より正である．よって矛盾する．

以上から $a \le y \le a - \dfrac{f(a)}{f'(a)}$ に $f(y) = 0$ なる $y$ がある．

\end{document}
